\chapter{Conclusion et remerciements}

\section{Synthèse du projet}

Ce projet de développement d'une application de \textbf{gestion de budget personnel} m'a permis de mettre en pratique les compétences acquises en alternance CDA dans un contexte concret et fonctionnel. L'architecture 3-tiers avec Next.js, Node.js/Express, PostgreSQL et MongoDB a démontré sa robustesse. Les objectifs métier ont été largement atteints : réduction de 28\,\% des dépassements budgétaires chez les testeurs et un taux de rétention de 83\,\% à J+30.

La démarche Agile m'a aidé à rester focalisé sur ce qui apporte vraiment de la valeur, et les bonnes pratiques de sécurité et de déploiement appliquées tout au long du projet garantissent la fiabilité de la solution livrée.

\begin{exemple}
\textbf{Chiffres clés du projet :}
\begin{center}
\begin{tabular}{|l|l|l|}
\hline
\textbf{Métrique} & \textbf{Valeur} & \textbf{Objectif} \\
\hline
Durée de développement & 5.5 mois & 6 mois \\
Couverture de code & 85{,}3\% & 80\% \\
Performance P95 & 342ms & 800ms \\
Vulnérabilités sécurité & 0 & 0 \\
Adoption utilisateurs & 78\% & 75\% \\
Temps de reporting & -42\% & -40\% \\
\hline
\end{tabular}
\end{center}

\textbf{Technologies maîtrisées :}
\begin{itemize}
    \item \textbf{Frontend :} Next.js 15, React 19, TypeScript
    \item \textbf{Backend :} Node.js, Express.js, pg-promise
    \item \textbf{Bases de données :} PostgreSQL 16 (pg-promise), MongoDB 7 (Mongoose)
    \item \textbf{DevOps :} Docker, GitHub Actions, SonarQube
    \item \textbf{Sécurité :} JWT, Bcrypt, OWASP Top 10
\end{itemize}
\end{exemple}




\section{Perspectives d'évolution}

La v2.0 de \textbf{My Smart Budget} s'articulera autour de trois axes qui me semblent vraiment utiles pour les utilisateurs : la synchronisation bancaire automatique (pour éliminer la saisie manuelle), une application mobile (pour consulter son budget partout), et des prévisions intelligentes basées sur l'historique. L'architecture actuelle a été pensée pour supporter ces évolutions sans refactoring majeur.

\begin{exemple}
\textbf{Roadmap technique v2.0 :}
\begin{verbatim}
Q1 2026: Qualité & Confort utilisateur
+-- Mode sombre (Dark Mode)
+-- Amélioration accessibilité mobile (polices, tap targets)
+-- Augmentation couverture tests composants Next.js (76% -> 90%)

Q2 2026: Intelligence artificielle
+-- POC catégorisation automatique (ML)
+-- Recommandations basées sur l'historique
+-- Détection de dépenses inhabituelles
+-- Prévisions de solde fin de mois

Q3 2026: Évolutions technologiques & Mobile
+-- Migration Node.js 22 LTS
+-- POC agrégation bancaire automatique (DSP2/Plaid)
+-- Application mobile PWA améliorée
+-- Monitoring Grafana dashboard dédié
\end{verbatim}

\textbf{Compétences à développer :}
\begin{itemize}
    \item \textbf{Architecture :} Microservices, Event-driven architecture
    \item \textbf{Cloud :} AWS/Azure, Kubernetes, Serverless
    \item \textbf{IA/ML :} TensorFlow, PyTorch, MLOps
    \item \textbf{Sécurité :} Zero Trust, DevSecOps
    \item \textbf{Leadership :} Architecture decision records, mentoring
\end{itemize}
\end{exemple}




\section{Remerciements}

Je tiens à remercier sincèrement toutes les personnes qui ont contribué à ce projet et à ma formation.

\begin{exemple}
\textbf{Remerciements :}
\begin{itemize}
    \item \textbf{Mon tuteur entreprise :} pour ses conseils techniques réguliers et sa disponibilité lors des phases de déploiement.
    \item \textbf{Les utilisateurs testeurs :} pour leurs retours honnêtes sur l'ergonomie et leur patience face aux premières versions.
    \item \textbf{Les formateurs CDA :} pour la transmission des fondamentaux techniques et méthodologiques qui ont structuré ce travail.
    \item \textbf{La communauté open source :} pour les outils (Next.js, Express.js, Docker, SonarQube, PostgreSQL) sans lesquels ce projet n'aurait pas été possible.
\end{itemize}
\end{exemple}




